\section{Adjoint kernels}
\label{sec:adjoints}

The reader is undoubtably familiar with the notion of adjoint
functors.  It is easy and natural to generalize this from the context
of the 2-category $\gCat$ of categories, functors and natural
transformations to the context of an arbitrary 2-category.  In this
section it is shown that every kernel, considered as a 1-morphism in
the 2-category $\gVar$, has both a left and right adjoint: this is a
consequence of Serre duality, and is closely related to the familiar
fact that every integral transform functor has both a left and right
adjoint functor.

Using these notions of left and right adjoints we define \emph{left
  partial trace maps}, and similarly \emph{right partial trace
  maps}. These can be viewed as partial versions of the Serre trace
map. This construction is very much the heart of the paper.

 \subsection{Adjunctions in 2-categories.}
The notion of an adjunction in a 2-category simultaneously generalizes
the notion of an adjunction between functors and the notion of a
duality between objects of a monoidal category.  As it is the former
that arises in the context of integral transforms, we will use that as
the motivation, but will come back to the latter below.


The most familiar definition of adjoint functors is as follows.  For
categories $\C$ and $\D$, an adjunction $\Psi\dashv \Phi$ between
functors $\Psi\colon \D\to \C$ and $\Phi\colon \C\to \D$ is the
specification of a natural isomorphism
\[t_{a,b}\colon \Hom_\C(\Psi(a),b)\xrightarrow{\sim}
\Hom_\D(a,\Phi(b))\]
for any $a\in \D$ and $b\in\C$.


It is well known (see \cite[page 91]{GelfandManin}) that this definition is
equivalent to an alternative definition of adjunction which consists
of the specification of unit and counit natural morphisms, namely
\[\eta\colon \Id_\D \Rightarrow \Phi\circ \Psi
\quad\text{and}\quad \epsilon\colon \Psi\circ \Phi\Rightarrow
\Id_\C, \]
such that the composite natural transformations
\[ \Psi \stackrel{\id_\Psi\circ\eta}\Longrightarrow
\Psi\circ\Phi\circ\Psi \stackrel{\epsilon\circ\Id_\Psi}\Longrightarrow
\Psi \quad\text{and}\quad \Phi
\stackrel{\eta\circ\Id_\Phi}\Longrightarrow \Phi\circ\Psi\circ\Phi
\stackrel{\Id_\Phi\circ\epsilon}\Longrightarrow \Phi
\]
are respectively the identity natural transformation on $\Psi$ and
the identity natural transformation on $\Phi$.


It is straight forward to translate between the two different
definitions of adjunction.  Given $\eta$ and $\epsilon$ as above, define
$t_{a,b}\colon\Hom_\C(\Psi(a),b)\xrightarrow{\sim}
\Hom_\D(a,\Phi(b))$ by $t_{a,b}(f):=\Phi(f)\circ \eta_a$.  The inverse
of $t_{a,b}$ is defined similarly using the counit $\epsilon$.
Conversely, to get the unit and counit from the natural isomorphism of
hom-sets, define $\eta_a:=t_{a,\Psi(a)}(\id_{\Psi(a)})$ and define
$\epsilon$ similarly.

The definition involving the unit and counit is stated purely in terms
of functors and natural transformations --- without mentioning objects
--- thus it generalizes immediately to arbitrary 2-categories.

\begin{Definition}
If $\gC$ is a 2-category, $X$ and $Y$ are objects of $\gC$, and
\[ \Phi\colon X \ra Y\quad \text{and}\quad \Psi\colon Y \ra X \]
are 1-morphisms, then an \emph{adjunction} between $\Phi$ and $\Psi$
consists of two 2-morphisms
\[ \eta\colon \Id_Y \Rightarrow \Phi\circ\Psi \quad\text{and}\quad
\epsilon\colon \Psi\circ\Phi \Rightarrow \Id_X, \]
such that
\[(\epsilon\circh\Id_\Psi)\circv(\Id_\Psi\circh\eta)=\Id_\Psi
\quad\text{and}\quad(\Id_\Phi\circh\epsilon)\circv(\eta\circh\Id_\Phi)=\Id_\Phi.\]
Given such an adjunction we write $\Psi\dashv\Phi$.
\end{Definition}


It is worth noting that this also generalizes the notion of duality in
a monoidal category, that is to say two objects are dual in a monoidal
category if and only if the corresponding 1-morphisms are adjoint in
the corresponding 2-category-with-one-object.  Indeed, taking this
point of view, May and Sigurdsson \cite{MaySigurdsson} refer to what
is here called adjunction as duality.


It is at this point that the utility of the string diagram notation
begins to be seen.  Given an adjunction $\Psi\dashv\Phi$ the counit
$\epsilon \colon \Psi\circ \Phi \Rightarrow\Id$ and the unit
$\eta\colon \Id\Rightarrow\Phi\circ \Psi$ can be denoted as follows:
\[\pic{NaiveCounit2Adjunction};\qquad \text{and}\qquad
\pic{NaiveUnit2Adjunction}.\]
However, adopting the convention of denoting the identity one-morphism
by omission, it is useful just to draw the unit and counit as a cup
and a cap respectively:
\[\pic{Counit2Adjunction}:=\pic{NaiveCounit2Adjunction};\qquad
\text{and}\qquad \pic{Unit2Adjunction}:= \pic{NaiveUnit2Adjunction}.\]
The relations become the satisfying
\[\pic{SnakePsi}=\pic{IdPsi}\quad\text{and}\quad\pic{SnakePhi}=\pic{IdPhi}.\]


Adjunctions in 2-categories, as defined above, do correspond to
isomorphisms of certain hom-sets but in a different way to the
classical notion of adjunction.  Namely, if $\Theta\colon Z\to Y$ and
$\Xi\colon Z\to X$ are two other 1-morphisms, then an adjunction
$\Psi\dashv \Phi$ as above gives an isomorphism
\begin{align*}
  \THom(\Psi\circ\Theta,\Xi)&\xrightarrow{\sim}\THom(\Theta,
  \Phi\circ \Xi)\\
  \pic{AdjunctionIso1}&\mapsto \pic{AdjunctionIso2}
\end{align*}
The inverse isomorphism uses the counit in the obvious way.

In a similar fashion, for $\widehat\Theta\colon Y\to Z$ and $\widehat\Xi\colon
X\to Z$ two 1-morphisms, one obtains an isomorphism
\[\THom(\widehat\Theta\circ\Phi,\widehat\Xi)\xrightarrow{\sim}\THom(\widehat\Theta,
\widehat\Xi\circ\Psi),\]
for which the reader is encouraged to draw the relevant pictures.  It
is worth noting that with respect to the previous isomorphism, $\Psi$
and $\Phi$ have swapped sides in all senses.


Adjunctions are unique up to a canonical isomorphism by the usual
argument.  This means that if $\Psi$ and $\Psi'$ are, say, both left
adjoint to $\Phi$, then there is a canonical isomorphism
$\Psi\stackrel{\sim}\Longrightarrow\Psi'$.  This is pictured below and
it is easy to check that this is an isomorphism.
\[\pic{AdjCanIso}\]

Adjunctions are natural in the sense that they are preserved by
2-func\-tors, so, for instance, given a pair of adjoint kernels in
$\gVar$, the corresponding integral transforms are adjoint functors.

 \subsection{Left and right adjoints of kernels.}
\label{subsec:leftrightadjoints}
\label{subsec:yphimphi}
In an arbitrary 2-category a given 1-morphism might or might not have
a left or a right adjoint, but in the 2-category $\gVar$ every
1-morphism, that is every kernel, has both a left and a right adjoint.
We will see below that for a kernel $\Phi\colon X\to Y$ there are
adjunctions
\[\Phi^\chk\circ\SerK_Y\adjoint\Phi\adjoint \SerK_X\circ
\Phi^\chk,\]
where $\Phi^\chk\colon Y\to X$ means the object $\sHom_{\D(X\times
  Y)}(\Phi,\cO_{X\times Y})$ considered as an object in $\D(Y\times
X)$.  This should be compared with the fact that if $M$ is an
$A$-$B$-bimodule then $M^\chk$ is naturally a $B$-$A$-bimodule.  We
shall see that the two adjunctions above are related in some very
useful ways.

\begin{Proposition}\label{prop:deltalstar}
  If $X$ is a space and $\Delta\colon X\to X\times X$ is the diagonal
  embedding then $\Delta_*\colon \D(X)\to \D(X\times X)$, the
  push-forward on derived categories, is a monoidal functor where
  $\D(X)$ has the usual monoidal tensor product $\otimes$ and
  $\D(X\times X)$ has the composition $\circ$ as the monoidal
  structure.
\end{Proposition}

\begin{Proof}
The proof is just an application of the projection formula.
\end{Proof}

\noindent This has the following immediate consequence.

\begin{Lemma}
  If $\cE$ and $\cF$ are dual as objects in $\D(X)$ then $\Delta_*\cE$
  and $\Delta_*\cF$ are both left and right adjoint to each other as
  1-morphisms in $\gVar$.
\end{Lemma}


In order not to hold-up the flow of the narrative, the proofs of the
remaining results from this section have been relegated to
Appendix~\ref{sec:duality}.

 We begin with some background on the
Serre kernel $\SerK_X$.
Recall from Section~\ref{sec:serre} that the anti-Serre kernel
$\SerK^{-1}_X$ is defined to be $\Delta_*\omega^{-1}_X[-\dim X]$, and
that the Serre kernel $\SerK_X$ is denoted by a dashed-dotted line,
while the anti-Serre kernel $\SerK^{-1}_X$ is denoted by a
dashed-dotted line with a horizontal bar.  The above proposition means
that we actually have maps
\[\pic{SerK1},\quad
\pic{SerK2},\quad
\pic{SerK3},\quad
\pic{SerK4}
\]
satisfying the following relations
\[\pic{SerK5}=\pic{SerK6},\quad
\pic{SerK7}=\pic{SerK8},\quad
\pic{SerK9}=\pic{SerK10},\quad
\]
and all obvious variations thereof.


%\label{subsec:yphimphi}
In the appendix it is shown that for a kernel $\Phi$ there are natural
morphisms $\rote_\Phi\colon \Phi\circ\SerK_X\circ\Phi^\chk\to \cO_\Delta$
and $\gamma_\Phi\colon \SerK_X^{-1}\to \Phi^\chk\circ \Phi$, denoted in the
following fashion, where the solid, upward oriented lines are labelled
with $\Phi$ and the solid, implicitly downward oriented lines are
labelled by $\Phi^\chk$:
\[\rote_\Phi:\pic{M23}\quad \text{and}\quad \gamma_\Phi:\pic{Y23c}.\]
The main property of $\rote_\Phi$ and $\gamma_\Phi$ is that if we define
$\epsilon_\Phi$, $\overline \epsilon_\Phi$, $\eta_\Phi$ and
$\overline\eta_\Phi$ via
\begin{gather*}
 \eta_\Phi: \pic{eta1}:=\pic{YbentRight},\qquad
\overline\eta_\Phi: \pic{eta2}:=\pic{YbentLeft},\\
\epsilon_\Phi:=\rote_{\Phi^\chk}=\pic{M14},\qquad
\overline\epsilon_\Phi:=\rote_{\Phi}=\pic{M23},
\end{gather*}
then these are the units and counits of adjunctions
\[\Phi^\chk\circ\SerK_Y\adjoint\Phi\adjoint \SerK_X\circ
\Phi^\chk.\]

\subsection{Partial traces.}
We can now define the important notion of partial traces.
\begin{Definition}
   For a kernel $\Phi\colon X\kerto Y$, and
  kernels $\Psi,\Theta\colon Z\kerto X$ define the \emph{left partial trace}
  \[ \THom(\Phi\circ\Theta,\SerK_Y\circ\Phi\circ\Psi) \to
  \THom(\Theta,\SerK_X\circ\Psi) \]
  as
  \[ \pic{UnpartialTraceLeft}\mapsto\pic{PartialTraceLeft}.\]
  Similarly we define a \emph{right partial trace}
  \[ \THom(\Theta'\circ\Phi,\Psi'\circ\Phi\circ\SerK_X) \to
  \THom(\Theta',\Psi'\circ\SerK_Y)\]
  as
  \[ \pic{UnpartialTraceRight}\mapsto\pic{PartialTraceRight}.\]
\end{Definition}
\medskip

\noindent
The following key result, proved in the appendix, says that taking partial
trace does not affect the Serre trace.

\begin{Theorem}
  For a kernel $\Phi\colon X\kerto Y$, a kernel $\Psi\colon Z\kerto X$
  and a kernel morphism
  $\alpha\in\Hom(\Phi\circ\cF,\SerK_Y\circ\Phi\circ\Psi\circ\SerK_Z)$
  then the left partial trace of $\alpha$ has the same Serre trace as
  $\alpha$, i.e., pictorially
  \[\Tr\left(\pic{ApPT1}\right)=\Tr\left(\pic{ApPT2}\right).\]
  The analogous result holds for the right partial trace.
\end{Theorem}

 \subsection{Adjunction as a 2-functor.}
\label{subsec:adj2functors}
As shown in Section~\ref{subsec:leftrightadjoints}, in the
2-cate\-gory $\gVar$ every 1-morphism, that is every kernel,
$\Phi\colon X\to Y$ has a right adjoint $\SerK_X\circ \Phi^\chk$.
This can be extended to a `right adjunction 2-functor' $\tau_R\colon
\gVar^\text{coop}\to \gVar$, where $\gVar^\text{coop}$ means the
contra-opposite 2-category of $\gVar$, which is the 2-category with
the same collections of objects, morphisms and 2-morphisms, but in
which the direction of the morphisms and the 2-morphisms are reversed.


Before defining $\tau_R$, however, it is perhaps useful to think of
the more familiar situation of a one-object 2-category with right
adjoints, i.e., a monoidal category with (right) duals.  So if $\gC$
is a monoidal category in which each object $a$ has a dual $a^\chk$
with evaluation map $\epsilon_a\colon a^\chk\otimes a\to \one$ and
coevaluation map $\eta_a\colon \one \to a\otimes a^\chk$, then for any
morphism $f\colon a\to b$ define $f^\chk\colon b^\chk\to a^\chk$ to be
the composite:
\[b^\chk\xrightarrow{\id\otimes \eta} b^\chk\otimes a\otimes
a^\chk\xrightarrow{\id\otimes f\otimes \id} b^\chk\otimes b\otimes
a^\chk\xrightarrow{\epsilon\otimes \id} a^\chk.\]
This gives rise to a functor $(-)^\chk\colon \gC^\text{op}\to \gC$.


Now return to the case of interest and define $\tau_R\colon
\gVar^\text{coop}\to \gVar$ as follows.  On spaces define
$\tau_R(X):=X$.  On a kernel $\Phi\colon X\to Y$ define
$\tau_R(\Phi):=\SerK_X\circ \Phi^\chk$.  Finally, on morphisms of
kernels define it as illustrated:
\[\tau_R\left(\pic{BeforetauRalpha}\right):=\pic{tauRalpha}.\]
It is a nice exercise for the reader to check that this is a
2-functor.


Clearly a left adjoint 2-functor $\tau_L\colon \gVar^\text{coop}\to
\gVar$ can be similarly created by defining it on a kernel $\Phi\colon
X\to Y$ by $\tau_L(\Phi):=\Phi^\chk\circ\SerK_Y$ and by defining it on
morphisms of kernels by
\[\tau_L\left(\pic{BeforetauRalpha}\right):=\pic{tauLalpha}.\]
